\begin{textsample}{NEG Dim 1 – 2009\_c5 – Score -1.00 – 2009\_c5/enade\_2009\_tec\_mark\_8016f58348.txt}  \label{ex:f1_neg_291}
Capítulo 5 Distribuição dos Conceitos 5.1 Panorama nacional da distribuição dos conceitos A Tabela 5.1 apresenta a quantidade e o percentual de cursos de TECNOLOGIA EM MARKETING participantes do ENADE/2009 que se situaram em \textbf{cada} uma das faixas de conceitos, distribuídos de acordo com a Região em que se localizam os cursos das instituições de educação superior, mantenedoras dos cursos responsáveis pela oferta.
O Gráfico 5.1 apresenta a mesma informação sob a forma de diagrama de caixa para uma melhor visualização da dispersão dos conceitos obtidos.
No diagrama de caixa, a caixa propriamente dita fica limitada pelos quartis.
A barra dentro da caixa em negrito representa a mediana.
Linhas verticais unem os quartis aos valores extremos.
Se estes valores extremos se configurarem como valores atípicos, estes aparecem como pontos isolados e as linhas terminam no maior ( ou menor ) valor ainda não considerado atípico.
Eventuais valores atípicos aparecem como pequenos círculos com a categoria administrativa correspondente.
Valores considerados muito atípicos foram representados por estrelas.
Um diagrama de caixa para poucos dados ou com dados muito concentrados pode ter alguns pontos coincidentes.
Por exemplo, se mais da metade dos cursos de uma Região receberam conceito 5, o valor extremo superior, o quartil superior e a mediana coincidirão.
Se o conjunto só contém um dado este apareceria denotado como a barra da mediana.
Se o conjunto apresenta todas as notas iguais, a notação é a mesma, somente uma barra horizontal.
Observando -se os dados da Tabela 5.1, vemos que os 222 cursos participantes na área de TECNOLOGIA EM MARKETING \textbf{estavam} distribuídos entre as Regiões como : 10 na Região Norte ( correspondendo a 4,5 \% dos cursos ), 26 na Região Nordeste ( correspondendo a 11,7 \% dos cursos ), 131 na Região Sudeste ( correspondendo a 59,0 \% dos cursos ), 39 na Região Sul ( correspondendo a 17,6 \% dos cursos ) e 16 na Região Centro-Oeste ( correspondendo a 7,2 \% dos cursos ).
Destes 222 cursos, 44 não foram avaliados : 2 na Região Norte, 7 na Região Nordeste, 27 na Região Sudeste, 5 na Região Sul e 3 na Região Centro-Oeste.
Percebe -se que, em termos nacionais, predominaram os cursos que receberam conceitos 2, 3 e 4, correspondendo respectivamente a 25,3 \%, 51,1 \% e a 16,9 \% do total nacional avaliado.
Pouco mais de 4 \% dos cursos foram classificados com conceito 1.
O conceito máximo foi alocado a 2,8 \% dos cursos com conceito.
Considerando -se separadamente as cinco Regiões, observou -se que, da Região Norte, 10 cursos participaram, dos quais 2 não foram avaliados.
Os 8 cursos avaliados ficaram alocados igualmente entre os conceitos 2 e 3.
Nenhum curso da Região obteve o conceito máximo e nenhum curso da Região obteve o conceito mínimo.
Como os cursos se concentraram em somente duas categorias contíguas, o diagrama de caixa se reduz à parte central.
Na Região Nordeste, 26 cursos participaram, dos quais, 7 não foram avaliados.
Entre os avaliados, o grupo modal entre os avaliados na Região, obteve conceito 3 ( 47,4 \% ).
Somente 1 curso desta Região obteve conceito mínimo, 5 cursos ( 26,3 \% dos avaliados na Região ) obtiveram conceito 2, 3 e conceito 4 ( 15,8 \% ) e 1 obteve o conceito máximo..
Na Região Sudeste, 131 cursos participaram, dos quais, cerca de 1/5 não foi avaliado ( 27 cursos ).
Dentre os avaliados, a maioria obteve conceito 3, equivalentes a 54,8 \% do total dos avaliados na Região.
O conceito máximo foi alcançado por 3 cursos ( 2,9 \% dos avaliados na Região ), outros 4 cursos obtiveram conceito mínimo, 19 obtiveram o conceito 4 ( 18,3 \% dos avaliados na Região ) e 21 ( 20,2 \% dos avaliados na Região ) receberam conceito 2.
Esta Região, juntamente com a Região Sul, foi a que apresentou cursos com a melhor distribuição de conceitos.
Na Região Sul, 5 dos 39 cursos participantes não foram avaliados.
Dentre os avaliados, equivalentes a 47,1 \% do total regional avaliado, classificaram -se com conceito 3.
O conceito máximo foi alcançado por somente 1 curso ( 2,9 \% dos avaliados ), 10 cursos ( 29,4 \% dos avaliados ) obtiveram conceito 2 e 7 ( 20,6 \% dos avaliados ) cursos ficaram com conceito 4.
Na Região Centro-Oeste, dos 16 cursos participantes 3 não foram avaliados.
Dentre os avaliados, 38,5 \% ( 5 cursos \textbf{cada} ) obtiveram conceito 2 ou 3, os conceitos modais.
O conceito máximo não foi alcançado por nenhum curso, 2 cursos ( 15,4 \% dos avaliados ) obtiveram conceito mínimo, 1 obteve o conceito 4.
Esta Região foi a que apresentou cursos com a pior distribuição de conceitos. 5.2 Conceitos por categoria administrativa e por Região A Tabela 5.2 apresenta a distribuição dos cursos participantes do ENADE/2009 de TECNOLOGIA EM MARKETING, por categoria administrativa, de acordo com os conceitos por eles alcançados, segundo Região.
O Gráfico 5.2 apresenta a mesma informação sob a forma de diagrama de caixa para uma melhor visualização da dispersão dos conceitos obtidos.
Somente eventuais valores atípicos aparecem com o identificador da Região na qual o curso foi habilitado.
De acordo com as informações apresentadas, a grande maioria, 98,6 \%, foi de cursos em instituições privadas e 1,4 \% dos cursos eram vinculados a instituições municipais.
Não foi avaliado 1 dos 3 cursos em instituições municipais e 43 dos 219 cursos em instituições privadas.
Os dois cursos avaliados em instituições municipais receberam conceito 2 e 3.
Os cursos avaliados nas instituições privadas tiveram como conceito modal o valor 3 com uma incidência de 51,1 \% dos avaliados, correspondendo a 90 cursos.
Somente 5 cursos ( 2,8 \% dos avaliados ) receberam o conceito máximo, 7 ( 4,0 \% dos avaliados ) receberam o conceito mínimo, 30 receberam conceito 2 e 44, o conceito 4.
Considerando -se a análise por Região observou -se que em todas as Regiões que as Instituições privadas \textbf{estavam} sempre representadas em todas as Regiões.
Somente nas Regiões Sul e Sudeste existem a presença de cursos em instituições municipais, nos outros os cursos \textbf{estão} exclusivamente em instituições privadas.
Na Região Norte, os 10 cursos participantes \textbf{estavam} concentrados em instituições privadas.
Destes, 2 não foram avaliados.
Os 8 cursos avaliados \textbf{estavam} distribuídos igualmente nos conceitos 2 e 3.
Nenhum curso da Região obteve nem o conceito máximo nem o conceito mínimo.
Na Região Nordeste, todos os 26 cursos participantes \textbf{estavam} alocados em instituições privadas.
Desses, 7 cursos não foram avaliados.
Entre os 19 avaliados, o conceito modal foi 3, com 47,4 \% dos avaliados na Região ( 9 cursos ).
Somente 1 curso desta Região obteve conceito mínimo e 1 outro, o conceito máximo ( 5,3 \% dos avaliados ).
Os demais cursos foram distribuídos nos conceitos 2 ( 5 cursos correspondendo a 26,3 \% dos avaliados na Região ) e no conceito 4 ( 3 cursos correspondendo a 15,8 \% dos avaliados ) Na Região Sudeste, dos 131 cursos participantes, 130 \textbf{estavam} alocados em instituições privadas.
Um curso \textbf{estava} alocado em instituição municipal, que recebeu conceito 3.
Dos 130 cursos participantes em instituições privadas na Região, um pouco mais de 1/5 não foi avaliado ( 27 cursos ).
Dentre os avaliados, a maioria obteve conceito 3, equivalentes a 54,4 \% do total dos avaliados em instituições privadas na Região.
O conceito máximo foi alcançado por 3 cursos ( 2,9 \% dos avaliados em instituições privadas na Região ), outros 4 cursos obtiveram conceito mínimo ( 20,4 \% ), 19 obtiveram o conceito 4 ( 18,4 \% dos avaliados ) e 21 ( 20,4 \% dos avaliados ) receberam conceito 2.
Na Região Sul, dos 39 cursos participantes, 37 \textbf{estavam} alocados em instituições privadas.
Os demais 2 cursos \textbf{estavam} em instituições municipais, dos quais apenas 1 foi avaliado e recebeu conceito 2.
Dos 37 cursos em instituições privadas, um pouco mais de 10 \% ( 4 cursos ) não foram avaliados.
Dentre os avaliados, 16 cursos, equivalentes a 48,5 \% do total regional avaliado em instituições privadas, classificaram -se com conceito 3.
O conceito máximo foi alcançado por somente 1 curso ( 3,0 \% ), 7 cursos ( 21,2 \% ) obtiveram conceito 4 e outros 9 ficaram com conceito 2 ( 27,3 \% ).
Na Região Centro-Oeste, todos os 16 cursos participantes \textbf{estavam} alocados em instituições privadas.
Dentre estes 16 cursos em instituições privadas, quase 1/5, 3 cursos, não foram avaliados.
Dentre os avaliados, 38,5 \% ( 5 cursos ) obtiveram conceitos 2 ou 3, conceitos modais.
O conceito máximo não foi alcançado por nenhum curso, 2 cursos ( 15,4 \% dos avaliados ) obteve conceito mínimo, 1 obteve o conceito 4 ( 7,7 \% dos avaliados ). 5.3 Conceitos por Organização Acadêmica e por Região Na Tabela 5.3 encontra -se a distribuição dos conceitos atribuídos aos cursos participantes do ENADE/2009 de TECNOLOGIA EM MARKETING, por Organização Acadêmica, segundo as Regiões.
O Gráfico 5.3 apresenta a mesma informação sob a forma de diagrama de caixa para uma melhor visualização da dispersão dos conceitos obtidos.
Somente eventuais valores atípicos apareceram com o identificador da Categoria Administrativa.
Os 222 cursos desta área \textbf{estão} principalmente alocados em Faculdades ( 56,3 \% correspondendo a 125 cursos ), depois em Universidades ( 27,0 \% correspondendo a 60 cursos ) e Centros Universitários ( 16,7 \% correspondendo a 37 cursos ).
Destes 222 cursos participantes, 44 não foram avaliados : 5 \textbf{estavam} alocados em Universidades, 7 em Centros Universitários e 32 em Faculdades.
O ordenamento das 3 categorias principais se repete somente nas Regiões Sudeste e Sul.
Na Região Norte e Centro-Oeste, os Centros Universitários \textbf{estão} em segundo lugar.
Nas 3 categorias principais, Universidades, Centros Universitários e Faculdades, o conceito modal foi 3 ( respectivamente 52,7 \%, 40,0 \% e 53,8 \% dos avaliados em \textbf{cada} categoria ).
Proporcionalmente poucos cursos nos 3 tipos de Organização Acadêmica participantes, Universidades, Centros Universitários e Faculdades, receberam conceito mínimo, respectivamente 5,5 \%, 6,7 \% e 2,2 \% dos avaliados em \textbf{cada} categoria.
Uma proporção quase da mesma \textbf{ordem} de grandeza atingiu o conceito máximo : 1,8 \% em Universidades e 4,3 \% em Faculdades.
Nenhum curso em Centros Universitários alcançou o conceito máximo.
Considerando -se separadamente as cinco Regiões brasileiras, observou -se que, na Região Norte, os 10 cursos \textbf{estavam} distribuídos : 1 em Universidade, 2 em Centros Universitários e 7 em Faculdades.
O único curso em Universidade recebeu conceito 3.
Dos 2 cursos em Centros Universitários, 1 não foi avaliado e 1 recebeu conceito 2.
Dos 7 cursos em Faculdades, 1 não foi avaliado e os outros 6 foram avaliados e distribuídos igualmente nos conceitos 2 e 3.
Na Região Nordeste, os 26 cursos \textbf{estavam} distribuídos nas 3 categorias principais : 3 em Universidades, 3 em Centros Universitários e 20 em Faculdades.
Dos 3 cursos em Universidades, todos foram avaliados e receberam conceitos 2, 3 e 5.
Dos 3 cursos em Centros Universitários, 2 foram avaliados e receberam conceitos 3 e 4.
O conceito modal entre os avaliados em Faculdades foi 3, com 7 cursos ( correspondendo a 50,0 \% dos avaliados na categoria na Região ).
Nenhum curso nesta categoria recebeu conceito mínimo e somente 1 recebeu o conceito máximo.
Na Região Sudeste os 131 cursos \textbf{estavam} distribuídos nas 3 categorias principais : 44 em Universidades, 25 em Centros Universitários e 62 em Faculdades.
Os 27 cursos não avaliados também \textbf{estavam} distribuídos nas 3 categorias principais : 3 em Universidades, 5 em Centros Universitários e 19 em Faculdades.
A maioria dos 41 cursos avaliados em Universidades na Região recebeu conceito 3 ( 24 cursos correspondendo a 58,5 \% dos avaliados ), categoria também modal para Centros Universitários ( com 8 dos 20 cursos avaliados ) e Faculdades ( com 25 dos 43 cursos avaliados ).
Centros Universitários não apresentaram cursos nem com o conceito mínimo, nem máximo.
Universidades e Faculdades apresentaram cursos com conceito mínimo ( respectivamente 4,9 \% e 4,7 \% dos avaliados em \textbf{cada} categoria ) e máximo ( respectivamente 2,4 \% e 4,7 \% dos avaliados em \textbf{cada} categoria ).
Na Região Sul, os 39 cursos \textbf{estavam} distribuídos somente nas 3 categorias mais frequentes : 10 em Universidades, 3 em Centros Universitários e 26 em Faculdades.
Os 5 cursos não avaliados \textbf{estavam} assim distribuídos : 2 em Universidades e 3 em Faculdades.
Em Universidades, os conceitos modais foram 3 e 4 ( com 3 cursos \textbf{cada} ).
O conceito modal para os demais tipos de Organização Acadêmica foi o conceito 3 para os Centros Universitários ( 66,7 \% dos avaliados ) e Faculdades ( 42,3 \% ).
O conceito máximo foi alcançado por somente 1 curso e em uma Faculdade.
O conceito mínimo não foi alocado na Região.
Na Região Centro-Oeste, os 16 cursos participantes \textbf{estavam} assim distribuídos : 2 em Universidades, 4 em Centros Universitários e 10 em Faculdades.
Os 3 cursos não avaliados na Região \textbf{estavam} todos em Faculdades.
Os dois cursos avaliados em Universidades receberam conceito 2.
Os 4 cursos avaliados em Centros Universitários receberam conceitos 1 ( 2 cursos ), 3 ( 1 curso ) e 4 ( 1 curso ).
Nas Faculdades, o conceito modal foi 3, com 4 cursos.
Os outros 3 cursos avaliados receberam conceito 2.
Nenhum curso em Faculdades recebeu nem o conceito mínimo, nem o máximo.

% matched lemmas: cada, estar, ordem
\end{textsample}
